\section{Introduction}
\label{sec:introduction}

Species histories are not always trees. Hybridization, introgression and
allopolyploidy move genetic material between diverging lineages, and the result
is a reticulate history that a tree cannot express
\cite{TODO_networks_review}. Inferring such histories is harder than inferring
trees for two reasons that compound: the space of networks is larger and less
well behaved than the space of trees, and the same gene-tree discordance that
hybridization produces is also produced by incomplete lineage sorting
\cite{TODO_ILS_review}. Distinguishing the two requires a model in which both
occur, which is what the multispecies network coalescent provides
\cite{TODO_MSNC}.

Over roughly fifteen years PhyloNet \cite{Than2008PhyloNet,
Wen2018InferringPhyloNet} assembled the methods this implies: parsimony,
maximum likelihood, pseudo-likelihood, and Bayesian inference, from gene trees,
from sequence alignments, and from biallelic markers, with extensions to
polyploid complexes. Its importance to the field is not in dispute and this
paper does not argue otherwise. A researcher today who wants to infer a
phylogenetic network from gene trees reaches for PhyloNet, and should.

What has changed is the surrounding computational environment. Three specific
frictions motivated this work.

First, a PhyloNet analysis is specified in a NEXUS file containing a command
block and run through the JVM. That design is well suited to running an analysis
and poorly suited to embedding one: a sweep over conditions means generating
files, and consuming the result means parsing text. Comparative and statistical
analysis in phylogenetics now largely happens in Python, alongside NumPy
\cite{Harris2020NumPy} and SciPy \cite{Virtanen2020SciPy}, and the tree-oriented
libraries that serve that environment -- DendroPy \cite{Sukumaran2010DendroPy},
ETE \cite{HuertaCepas2016ETE3}, Biopython \cite{Cock2009Biopython} -- have no
network counterpart.

Second, a method developer starting a new network method starts from very little.
A network representation, traversal and cluster machinery, topology moves, a
search driver, parameter optimisation, input parsing and convergence diagnostics
are all prerequisites to testing a new scoring criterion, and none of them are
what the new method is about. The predictable outcome is that each new method
ships as a separate program with its own representation, and that comparing two
methods means reconciling two file formats.

Third, the organisation of a method set as a list of command names hides its
structure. Most of those commands vary along the same three dimensions -- the
data, the assumed generative process, the optimality criterion -- and a list
does not say which combinations exist, which are unimplemented, and which are
not meaningful at all. A user cannot tell the difference between a feature that
is missing and a request that does not make sense.

\phynetpy\ addresses these. It is a Python library in which the network is a
reusable data structure, inference is two verbs over three typed axes, and the
set of available methods is queryable at runtime. Its inference methods are, for
the most part, reimplementations of PhyloNet's, which is deliberate: the point is
not new statistics but a foundation on which new statistics can be built and
against which existing results can be checked.

\paragraph{Contributions.}
\begin{itemize}
  \item A network representation and analysis layer usable independently of any
    inference method: a rooted DAG with ploidy, typed branch-length units, nine
    dissimilarity measures, blob and cluster decomposition, and MUL-tree
    conversion (Sections~\ref{sec:architecture} and~\ref{sec:functionality}).
  \item An inference interface of two verbs over the product of data, model and
    criterion, with dispatch through a registry that distinguishes undefined,
    unimplemented and ill-typed requests, and that serves as the extension point
    for contributed methods (Section~\ref{sec:architecture}).
  \item Seven implemented inference methods, six corresponding to published
    PhyloNet analyses, with the gaps relative to PhyloNet stated explicitly
    (Section~\ref{sec:functionality}).
  \item Numerical cross-validation against PhyloNet's own implementations of the
    multispecies-network-coalescent density and the Felsenstein likelihood,
    across 14 adversarial cases (Section~\ref{sec:validation}).
  \item A replicated head-to-head runtime and accuracy comparison on identical
    input, reporting a runtime advantage and an accuracy deficit, with the
    accuracy result left open rather than explained away
    (Section~\ref{sec:validation}).
  \item A migration account for existing PhyloNet users, including what does not
    carry over (Section~\ref{sec:comparison}).
\end{itemize}
